\scanpage{104}
peut se définir comme à \uline{valeurs dans} $\mathcal A$ (et non seulement comme à val. dans $(\mathcal A_1\otimes\mathcal A_2)^{\unsure{\natural}}$, alors qu'on ne sait pas si $\mathcal A\longrightarrow(\mathcal A_1\otimes\mathcal A_2)^{\unsure{\natural}}$ est \inserted{seulement} injectif, ni d'ailleurs surjectif).

On doit pouvoir définir de même une norme, définie universellement par la condition que son image dans $(\mathcal A_1\otimes\mathcal A_2)^{\unsure{\natural}}$ soit $\xi.\bar\xi$.
Choisissons les bases
\[
(1,\unsure{U_1})\text{ de }\mathcal A_1,\qquad (1,U_2)\text{ de }\mathcal A_2
\]
donc
\[
\mathcal A_1=\mathcal O[U_1]/(U_1^2+b_1U_1+c_1),\qquad
\mathcal A_2=\mathcal O[U_2]/(U_2^2+b_2U_2+c_2),
\]
on aura
\[
\left\{\begin{aligned}
\mathcal A_1\otimes\mathcal A_2&=\mathcal O[U_1,U_2]/(U_1^2+b_1U_1+c_1,\ U_2^2+b_2U_2+c_2)\\
\mathcal A&\simeq\mathcal O[U]/(U^2+bU+c)
\end{aligned}\right.
\]
avec
\[
\left\{\begin{aligned}
b&=b_1b_2\\
c&=b_1^2c_2+b_2^2c_1-4c_1c_2
\end{aligned}\right.
\]
et $\mathcal A\longrightarrow\mathcal A_1\otimes\mathcal A_2$ est donné par
\[
U\longmapsto2U_1U_2+b_1U_2+b_2U_1
\]
Explicitons
\[
\mathcal A_1\otimes\mathcal A_2\xrightarrow[\inserted{x_1\otimes x_2\mapsto x_1*x_2}]{\tau}\mathcal A
\]
il suffit de trouver les images de $1,U_1\otimes1,1\otimes U_2,U_1\unsure{\otimes U_2}$.
\margin{c'est donc ce qui pour le \unsure{calcul} d'une \unsure{Tr.} de $\mathcal A_1\otimes\mathcal A_2$ \unsure{en disposant des}}
On aura bien sûr
\[
\deleted{\unsure{x_1*}}\quad x_1*e_1=\operatorname{Tr}_1(x_1).1_{\mathcal A}
\qquad e_1*x_2=\operatorname{Tr}_2(x_2).1_{\mathcal A}
\]
\margin{Ce qui définit $*$ sur $\mathcal A_1\unsure{\oplus|\amalg}\mathcal A_2\subset\mathcal A_1\otimes\mathcal A_2$.}
\mspage{103}
